% Options for packages loaded elsewhere
\PassOptionsToPackage{unicode}{hyperref}
\PassOptionsToPackage{hyphens}{url}
\documentclass[
  british,
  a4paper,
]{article}
\usepackage{xcolor}
\usepackage[margin=1in]{geometry}
\usepackage{amsmath,amssymb}
\setcounter{secnumdepth}{-\maxdimen} % remove section numbering
\usepackage{iftex}
\ifPDFTeX
  \usepackage[T1]{fontenc}
  \usepackage[utf8]{inputenc}
  \usepackage{textcomp} % provide euro and other symbols
\else % if luatex or xetex
  \usepackage{unicode-math} % this also loads fontspec
  \defaultfontfeatures{Scale=MatchLowercase}
  \defaultfontfeatures[\rmfamily]{Ligatures=TeX,Scale=1}
\fi
\usepackage{lmodern}
\ifPDFTeX\else
  % xetex/luatex font selection
  \setmainfont[]{Noto Serif}
  \setmonofont[]{Noto Sans Mono}
\fi
% Use upquote if available, for straight quotes in verbatim environments
\IfFileExists{upquote.sty}{\usepackage{upquote}}{}
\IfFileExists{microtype.sty}{% use microtype if available
  \usepackage[]{microtype}
  \UseMicrotypeSet[protrusion]{basicmath} % disable protrusion for tt fonts
}{}
\makeatletter
\@ifundefined{KOMAClassName}{% if non-KOMA class
  \IfFileExists{parskip.sty}{%
    \usepackage{parskip}
  }{% else
    \setlength{\parindent}{0pt}
    \setlength{\parskip}{6pt plus 2pt minus 1pt}}
}{% if KOMA class
  \KOMAoptions{parskip=half}}
\makeatother
% definitions for citeproc citations
\NewDocumentCommand\citeproctext{}{}
\NewDocumentCommand\citeproc{mm}{%
  \begingroup\def\citeproctext{#2}\cite{#1}\endgroup}
\makeatletter
 % allow citations to break across lines
 \let\@cite@ofmt\@firstofone
 % avoid brackets around text for \cite:
 \def\@biblabel#1{}
 \def\@cite#1#2{{#1\if@tempswa , #2\fi}}
\makeatother
\newlength{\cslhangindent}
\setlength{\cslhangindent}{1.5em}
\newlength{\csllabelwidth}
\setlength{\csllabelwidth}{3em}
\newenvironment{CSLReferences}[2] % #1 hanging-indent, #2 entry-spacing
 {\begin{list}{}{%
  \setlength{\itemindent}{0pt}
  \setlength{\leftmargin}{0pt}
  \setlength{\parsep}{0pt}
  % turn on hanging indent if param 1 is 1
  \ifodd #1
   \setlength{\leftmargin}{\cslhangindent}
   \setlength{\itemindent}{-1\cslhangindent}
  \fi
  % set entry spacing
  \setlength{\itemsep}{#2\baselineskip}}}
 {\end{list}}
\usepackage{calc}
\newcommand{\CSLBlock}[1]{\hfill\break\parbox[t]{\linewidth}{\strut\ignorespaces#1\strut}}
\newcommand{\CSLLeftMargin}[1]{\parbox[t]{\csllabelwidth}{\strut#1\strut}}
\newcommand{\CSLRightInline}[1]{\parbox[t]{\linewidth - \csllabelwidth}{\strut#1\strut}}
\newcommand{\CSLIndent}[1]{\hspace{\cslhangindent}#1}
\ifLuaTeX
\usepackage[bidi=basic,shorthands=off]{babel}
\else
\usepackage[bidi=default,shorthands=off]{babel}
\fi
\ifPDFTeX
\else
\babelfont{rm}[]{Noto Serif}
\fi
\ifLuaTeX
  \usepackage{selnolig} % disable illegal ligatures
\fi
\setlength{\emergencystretch}{3em} % prevent overfull lines
\providecommand{\tightlist}{%
  \setlength{\itemsep}{0pt}\setlength{\parskip}{0pt}}
\usepackage{array}
\usepackage{tabularx}
\usepackage{bookmark}
\IfFileExists{xurl.sty}{\usepackage{xurl}}{} % add URL line breaks if available
\urlstyle{same}
\hypersetup{
  pdflang={en-GB},
  hidelinks,
  pdfcreator={LaTeX via pandoc}}

\author{}
\date{}

\begin{document}

\section{Regular-entry prose pilot
02}\label{regular-entry-prose-pilot-02}

\emph{Assembly-only pilot for newly drafted compact regular-entry
commentary. The trace rendering is reused, but the prose comes from a
separate book-prose layer rather than from the model-entry report
sections.}

\subsection{bake --- OE bacan}\label{bake-oe-bacan}

Derivation: \emph{*bákaną} \textgreater{} \emph{bacan} (regular).

\begingroup
\setlength{\fboxsep}{6pt}

\noindent\fbox{%
\begin{minipage}{0.97\linewidth}
\small
\begin{tabularx}{\linewidth}{@{}>{\raggedright\arraybackslash}p{0.300\linewidth}>{\centering\arraybackslash}X>{\raggedright\arraybackslash}p{0.540\linewidth}@{}}
\begin{minipage}[t]{\linewidth}
\centering\textbf{Earlier Germanic changes}\par
\vspace{0.35em}
\raggedright
\centering\textbf{West Germanic}\par
\raggedright
\vspace{0.2em}
\raggedright [no change]\par
\vspace{0.6em}
\centering\textbf{Northwest Germanic}\par
\raggedright
\vspace{0.2em}
\raggedright [no change]\par
\end{minipage}
&
&
\begin{minipage}[t]{\linewidth}
\centering\textbf{Old English changes}\par
\vspace{0.35em}
\raggedright
\begin{tabular}{@{}>{\raggedright\arraybackslash}p{0.74\linewidth}@{\hspace{0.55em}}>{\raggedright\arraybackslash}p{0.16\linewidth}@{\hspace{0.25em}}}
Anglo Frisian Brightening & \emph{*bækaną} \\
OE A Restoration & \emph{*bakaną} \\
OE Heavy Syllable Nasal Apocope & \emph{*bakan} \\
OE Secondary Nasalization & \emph{*bakąn} \\
OE Weak Tail Reduction & \emph{*bakan} \\
\end{tabular}
\end{minipage}
\\
\end{tabularx}
\end{minipage}%
} \endgroup

Orel's \emph{*bakanan} and Campbell's standard examples of A-restoration
both point to the ordinary weak verb continued by Old English
\emph{bacan} (\citeproc{ref-Orel2003}{Orel 2003};
\citeproc{ref-Campbell1959}{Campbell 1959, 61};
\citeproc{ref-RingeTaylor2014}{Ringe and Taylor 2014, 270}). The
dictionary evidence is uncomplicated: \emph{bacan} is the attested
infinitive headword in Old English
(\citeproc{ref-BosworthToller1898}{Bosworth and Toller 1898, 72};
\citeproc{ref-ClarkHall1960}{Clark Hall 1960, 19}). What matters here is
the regular phonology. After Anglo-Frisian brightening, A-restoration
before single \emph{k} returns the stem vowel to \emph{a}, and the later
weak-verb tail reductions produce the expected infinitive.

\subsection{beech --- OE bōc}\label{beech-oe-bux14dc}

Derivation: \emph{*bōkō} \textgreater{} \emph{bōc} (regular).

\begingroup
\setlength{\fboxsep}{6pt}

\noindent\fbox{%
\begin{minipage}{0.97\linewidth}
\small
\begin{tabularx}{\linewidth}{@{}>{\raggedright\arraybackslash}p{0.460\linewidth}>{\centering\arraybackslash}X>{\raggedright\arraybackslash}p{0.460\linewidth}@{}}
\begin{minipage}[t]{\linewidth}
\centering\textbf{Earlier Germanic changes}\par
\vspace{0.35em}
\raggedright
\centering\textbf{West Germanic}\par
\raggedright
\vspace{0.2em}
\raggedright [no change]\par
\vspace{0.6em}
\centering\textbf{Northwest Germanic}\par
\raggedright
\vspace{0.2em}
\begin{tabular}{@{}>{\raggedright\arraybackslash}p{0.74\linewidth}@{\hspace{0.55em}}>{\raggedright\arraybackslash}p{0.16\linewidth}@{\hspace{0.25em}}}
\mbox{NWGmc Final Long O Raising} & \emph{*bōku} \\
\end{tabular}
\end{minipage}
&
&
\begin{minipage}[t]{\linewidth}
\centering\textbf{Old English changes}\par
\vspace{0.35em}
\raggedright
\begin{tabular}{@{}>{\raggedright\arraybackslash}p{0.74\linewidth}@{\hspace{0.55em}}>{\raggedright\arraybackslash}p{0.16\linewidth}@{\hspace{0.25em}}}
\mbox{OE High Vowel Apocope} & \emph{*bōk} \\
\end{tabular}
\end{minipage}
\\
\end{tabularx}
\end{minipage}%
} \endgroup

Kroonen's \emph{*bōk(j)ō-} already distinguishes nominative \emph{boc}
from oblique \emph{bēce}, so the relevant comparator here is the
nominative-singular line represented as \emph{bōc}
(\citeproc{ref-Kroonen2013}{Kroonen 2013, 81}). With \emph{*bōkō} as
input, the derivation is as compact as the trace suggests: final long
\emph{ō} raises in Northwest Germanic, then high-vowel apocope leaves
\emph{bōc}. The prose therefore only needs one reminder: the oblique
\emph{bēċe} belongs to the same paradigm, but not to the form chosen for
direct comparison.

\subsection{bier --- OE bǣr}\label{bier-oe-bux1e3r}

Derivation: \emph{*bḗrō} \textgreater{} \emph{bǣr} (regular).

\begingroup
\setlength{\fboxsep}{6pt}

\noindent\fbox{%
\begin{minipage}{0.97\linewidth}
\small
\begin{tabularx}{\linewidth}{@{}>{\raggedright\arraybackslash}p{0.460\linewidth}>{\centering\arraybackslash}X>{\raggedright\arraybackslash}p{0.460\linewidth}@{}}
\begin{minipage}[t]{\linewidth}
\centering\textbf{Earlier Germanic changes}\par
\vspace{0.35em}
\raggedright
\centering\textbf{West Germanic}\par
\raggedright
\vspace{0.2em}
\raggedright [no change]\par
\vspace{0.6em}
\centering\textbf{Northwest Germanic}\par
\raggedright
\vspace{0.2em}
\begin{tabular}{@{}>{\raggedright\arraybackslash}p{0.74\linewidth}@{\hspace{0.55em}}>{\raggedright\arraybackslash}p{0.16\linewidth}@{\hspace{0.25em}}}
\mbox{NWGmc Final Long O Raising} & \emph{*bḗru} \\
\mbox{NWGmc Long E Lowering} & \emph{*bǣru} \\
\end{tabular}
\end{minipage}
&
&
\begin{minipage}[t]{\linewidth}
\centering\textbf{Old English changes}\par
\vspace{0.35em}
\raggedright
\begin{tabular}{@{}>{\raggedright\arraybackslash}p{0.74\linewidth}@{\hspace{0.55em}}>{\raggedright\arraybackslash}p{0.16\linewidth}@{\hspace{0.25em}}}
\mbox{OE High Vowel Apocope} & \emph{*bǣr} \\
\end{tabular}
\end{minipage}
\\
\end{tabularx}
\end{minipage}%
} \endgroup

Kroonen's \emph{*bērō-} and the Old English evidence \emph{bar},
\emph{bær} point to an uncomplicated feminine noun continued here as
normalized \emph{bǣr} (\citeproc{ref-Kroonen2013}{Kroonen 2013, 717};
\citeproc{ref-ClarkHall1960}{Clark Hall 1960, 32};
\citeproc{ref-BosworthToller1898}{Bosworth and Toller 1898, 73}). The
only real editorial issue is spelling: the dictionaries transmit
\emph{bær} or \emph{bar}, while the book form \emph{bǣr} simply marks
the long vowel explicitly. Once that normalization is stated, the
derivation remains regular: final long \emph{ō} raises, long \emph{ē}
lowers to \emph{ǣ}, and apocope yields the expected Old English form.

\subsection{both --- OE bū}\label{both-oe-bux16b}

Derivation: \emph{*bō} \textgreater{} \emph{bū} (regular).

\begingroup
\setlength{\fboxsep}{6pt}

\noindent\fbox{%
\begin{minipage}{0.97\linewidth}
\small
\begin{tabularx}{\linewidth}{@{}>{\raggedright\arraybackslash}p{0.540\linewidth}>{\centering\arraybackslash}X>{\raggedright\arraybackslash}p{0.300\linewidth}@{}}
\begin{minipage}[t]{\linewidth}
\centering\textbf{Earlier Germanic changes}\par
\vspace{0.35em}
\raggedright
\centering\textbf{West Germanic}\par
\raggedright
\vspace{0.2em}
\raggedright [no change]\par
\vspace{0.6em}
\centering\textbf{Northwest Germanic}\par
\raggedright
\vspace{0.2em}
\begin{tabular}{@{}>{\raggedright\arraybackslash}p{0.74\linewidth}@{\hspace{0.55em}}>{\raggedright\arraybackslash}p{0.16\linewidth}@{\hspace{0.25em}}}
NWGmc Stressed Monosyllable O Raising & \emph{*bū} \\
\end{tabular}
\end{minipage}
&
&
\begin{minipage}[t]{\linewidth}
\centering\textbf{Old English changes}\par
\vspace{0.35em}
\raggedright
\raggedright [no change]\par
\end{minipage}
\\
\end{tabularx}
\end{minipage}%
} \endgroup

The regular item here is not the better-known masculine dual
\emph{bēġen} but the inherited neuter dual \emph{bū}. Kroonen's paradigm
under \emph{*ba-} includes precisely that unextended \emph{*bō} form
alongside the wider dual set, and the Old English grammars likewise
preserve neuter \emph{bū} beside masculine \emph{bēġen} and feminine
\emph{bā} (\citeproc{ref-Kroonen2013}{Kroonen 2013, 47};
\citeproc{ref-SieversBrunner1965}{Sievers 1965, 324};
\citeproc{ref-Campbell1959}{Campbell 1959, 295};
\citeproc{ref-Fulk2018}{Fulk 2018, 223}). That makes \emph{bū} the
cleanest regular comparator for this entry: it is attested Old English,
and it does not depend on the more debated extended formation behind
\emph{bēġen}. Campbell and Brunner both treat final stressed \emph{ō
\textgreater{} ū} as an ordinary West Germanic-to-Old-English
development, with \emph{bū} among the standard examples
(\citeproc{ref-Campbell1959}{Campbell 1959, 122};
\citeproc{ref-SieversBrunner1965}{Sievers 1965, 69}).

Comparison note. The historical discussion of \emph{bēġen} and the
continental forms behind Modern English \emph{both} belongs to a related
but different problem. Orel still cites an older \emph{*bō-j-}
explanation, while Fulk reports Seebold's preference for \emph{*bō-þ-}
(\citeproc{ref-Orel2003}{Orel 2003, 45}; \citeproc{ref-Fulk2018}{Fulk
2018, 223}). Those formations matter for the wider history of the
numeral, but they do not displace the regular comparison \emph{*bō
\textgreater{} bū} used here.

\subsection{bow --- OE bīeġan}\label{bow-oe-bux12beux121an}

Derivation: \emph{*báugijaną} \textgreater{} \emph{bīeġan} (regular).

\begingroup
\setlength{\fboxsep}{6pt}

\noindent\fbox{%
\begin{minipage}{0.97\linewidth}
\footnotesize
\begin{tabularx}{\linewidth}{@{}>{\raggedright\arraybackslash}p{0.280\linewidth}>{\centering\arraybackslash}X>{\raggedright\arraybackslash}p{0.560\linewidth}@{}}
\begin{minipage}[t]{\linewidth}
\centering\textbf{Earlier Germanic changes}\par
\vspace{0.35em}
\raggedright
\centering\textbf{West Germanic}\par
\raggedright
\vspace{0.2em}
\raggedright [no change]\par
\vspace{0.6em}
\centering\textbf{Northwest Germanic}\par
\raggedright
\vspace{0.2em}
\raggedright [no change]\par
\end{minipage}
&
&
\begin{minipage}[t]{\linewidth}
\centering\textbf{Old English changes}\par
\vspace{0.35em}
\raggedright
\begin{tabular}{@{}>{\raggedright\arraybackslash}p{0.62\linewidth}@{\hspace{0.55em}}>{\raggedright\arraybackslash}p{0.28\linewidth}@{\hspace{0.25em}}}
OE Au Fronting & \emph{*báeugijaną} \\
OE Diphthong Leveling & \emph{*bēagijaną} \\
OE Heavy Syllable Nasal Apocope & \emph{*bēagijan} \\
OE Secondary Nasalization & \emph{*bēagijąn} \\
Sievers Law Syncope & \emph{*bēagjąn} \\
OE Velar Palatalization & \emph{*bēaʤjąn} \\
OE I Umlaut & \emph{*bīeʤjąn} \\
OE Weak Tail Reduction & \emph{*bīeʤjan} \\
OE J Loss After Heavy & \emph{*bīeʤan} \\
\end{tabular}
\end{minipage}
\\
\end{tabularx}
\end{minipage}%
} \endgroup

Kroonen's weak causative \emph{*baugjan-} and Ringe and Taylor's fuller
chain to West Saxon \emph{biegan} place this verb firmly within the
ordinary bend-family rather than among analogical repair cases
(\citeproc{ref-Kroonen2013}{Kroonen 2013, 75};
\citeproc{ref-RingeTaylor2014}{Ringe and Taylor 2014, 94}). Old English
dictionary evidence is straightforward, and the spelling \emph{bīeġan}
here is simply a normalized representation of the attested weak verb
(\citeproc{ref-ClarkHall1960}{Clark Hall 1960, 34};
\citeproc{ref-BosworthToller1898}{Bosworth and Toller 1898, 102}). The
regularity lies in the sequence after pre-Old-English \emph{*bēagjan}:
palatalization of \emph{gj} and i-mutation produce the West Saxon
infinitive, so the prose only needs to identify the causative lexeme and
confirm that the output belongs to the expected verbal line.

\subsection{fare --- OE faran}\label{fare-oe-faran}

Derivation: \emph{*fáraną} \textgreater{} \emph{faran} (regular).

\begingroup
\setlength{\fboxsep}{6pt}

\noindent\fbox{%
\begin{minipage}{0.97\linewidth}
\small
\begin{tabularx}{\linewidth}{@{}>{\raggedright\arraybackslash}p{0.300\linewidth}>{\centering\arraybackslash}X>{\raggedright\arraybackslash}p{0.540\linewidth}@{}}
\begin{minipage}[t]{\linewidth}
\centering\textbf{Earlier Germanic changes}\par
\vspace{0.35em}
\raggedright
\centering\textbf{West Germanic}\par
\raggedright
\vspace{0.2em}
\raggedright [no change]\par
\vspace{0.6em}
\centering\textbf{Northwest Germanic}\par
\raggedright
\vspace{0.2em}
\raggedright [no change]\par
\end{minipage}
&
&
\begin{minipage}[t]{\linewidth}
\centering\textbf{Old English changes}\par
\vspace{0.35em}
\raggedright
\begin{tabular}{@{}>{\raggedright\arraybackslash}p{0.74\linewidth}@{\hspace{0.55em}}>{\raggedright\arraybackslash}p{0.16\linewidth}@{\hspace{0.25em}}}
Anglo Frisian Brightening & \emph{*færaną} \\
OE A Restoration & \emph{*faraną} \\
OE Heavy Syllable Nasal Apocope & \emph{*faran} \\
OE Secondary Nasalization & \emph{*farąn} \\
OE Weak Tail Reduction & \emph{*faran} \\
\end{tabular}
\end{minipage}
\\
\end{tabularx}
\end{minipage}%
} \endgroup

Kroonen and Orel both treat \emph{*faran} as the inherited strong verb,
and Campbell uses it as a standard example of A-restoration
(\citeproc{ref-Kroonen2013}{Kroonen 2013, 149};
\citeproc{ref-Orel2003}{Orel 2003, 132};
\citeproc{ref-Campbell1959}{Campbell 1959, 61}). The Old English
evidence is equally simple: dictionaries distinguish the strong
infinitive \emph{faran} from weak \emph{færan}, while forms such as
\emph{fære} and \emph{færð} belong to the present tense, not to the
citation infinitive (\citeproc{ref-ClarkHall1960}{Clark Hall 1960, 113};
\citeproc{ref-BosworthToller1898}{Bosworth and Toller 1898, 250}). That
is enough to explain why the regular comparison stays with \emph{faran}.
After Anglo-Frisian brightening, A-restoration before single \emph{r}
returns the stem vowel to \emph{a}, and the later tail reductions yield
the expected infinitive.

\subsection{guest --- OE ġiest}\label{guest-oe-ux121iest}

Derivation: \emph{*gástiz} \textgreater{} \emph{ġiest} (regular).

\begingroup
\setlength{\fboxsep}{6pt}

\noindent\fbox{%
\begin{minipage}{0.97\linewidth}
\small
\begin{tabularx}{\linewidth}{@{}>{\raggedright\arraybackslash}p{0.460\linewidth}>{\centering\arraybackslash}X>{\raggedright\arraybackslash}p{0.460\linewidth}@{}}
\begin{minipage}[t]{\linewidth}
\centering\textbf{Earlier Germanic changes}\par
\vspace{0.35em}
\raggedright
\centering\textbf{West Germanic}\par
\raggedright
\vspace{0.2em}
\raggedright [no change]\par
\vspace{0.6em}
\centering\textbf{Northwest Germanic}\par
\raggedright
\vspace{0.2em}
\begin{tabular}{@{}>{\raggedright\arraybackslash}p{0.74\linewidth}@{\hspace{0.55em}}>{\raggedright\arraybackslash}p{0.16\linewidth}@{\hspace{0.25em}}}
\mbox{PGmc Final Z Deletion} & \emph{*gásti} \\
\end{tabular}
\end{minipage}
&
&
\begin{minipage}[t]{\linewidth}
\centering\textbf{Old English changes}\par
\vspace{0.35em}
\raggedright
\begin{tabular}{@{}>{\raggedright\arraybackslash}p{0.74\linewidth}@{\hspace{0.55em}}>{\raggedright\arraybackslash}p{0.16\linewidth}@{\hspace{0.25em}}}
Anglo Frisian Brightening & \emph{*gæsti} \\
OE Velar Palatalization & \emph{*ʤæsti} \\
OE I Umlaut & \emph{*ʤesti} \\
OE Ws Palatal Diphthongization & \emph{*ʤiesti} \\
\mbox{OE High Vowel Apocope} & \emph{*ʤiest} \\
\end{tabular}
\end{minipage}
\\
\end{tabularx}
\end{minipage}%
} \endgroup

Campbell and Ringe and Taylor treat this as an ordinary i-stem whose
West Saxon development produces \emph{ġiest}, while non-West-Saxon
evidence preserves forms of the \emph{gest} type
(\citeproc{ref-Campbell1959}{Campbell 1959, 91};
\citeproc{ref-RingeTaylor2014}{Ringe and Taylor 2014, 296}).
Bosworth-Toller and Clark Hall indeed record a family of spellings such
as \emph{gist}, \emph{gest}, \emph{giest}, and \emph{gyst}, so the issue
is not whether the West Saxon form exists, but which member of that
attested set the book should foreground
(\citeproc{ref-BosworthToller1898}{Bosworth and Toller 1898, 387};
\citeproc{ref-ClarkHall1960}{Clark Hall 1960, 149}). The derivation
itself is regular: brightening and i-mutation reshape the stem, and West
Saxon palatal diphthongization produces \emph{ie}.

Dialect note. Anglian \emph{gest} remains genuine Old English evidence,
but it is a parallel dialectal comparator rather than a correction to
West Saxon \emph{ġiest} (\citeproc{ref-RingeTaylor2014}{Ringe and Taylor
2014, 296}; \citeproc{ref-BosworthToller1898}{Bosworth and Toller 1898,
387}).

\subsection{learn --- OE liornian}\label{learn-oe-liornian}

Derivation: \emph{*líznōjaną} \textgreater{} \emph{liornian} (regular).

\begingroup
\setlength{\fboxsep}{6pt}

\noindent\fbox{%
\begin{minipage}{0.97\linewidth}
\footnotesize
\begin{tabularx}{\linewidth}{@{}>{\raggedright\arraybackslash}p{0.280\linewidth}>{\centering\arraybackslash}X>{\raggedright\arraybackslash}p{0.560\linewidth}@{}}
\begin{minipage}[t]{\linewidth}
\centering\textbf{Earlier Germanic changes}\par
\vspace{0.35em}
\raggedright
\centering\textbf{West Germanic}\par
\raggedright
\vspace{0.2em}
\raggedright [no change]\par
\vspace{0.6em}
\centering\textbf{Northwest Germanic}\par
\raggedright
\vspace{0.2em}
\raggedright [no change]\par
\end{minipage}
&
&
\begin{minipage}[t]{\linewidth}
\centering\textbf{Old English changes}\par
\vspace{0.35em}
\raggedright
\begin{tabular}{@{}>{\raggedright\arraybackslash}p{0.62\linewidth}@{\hspace{0.55em}}>{\raggedright\arraybackslash}p{0.28\linewidth}@{\hspace{0.25em}}}
OE Breaking & \emph{*líornōjaną} \\
OE Heavy Syllable Nasal Apocope & \emph{*líornōjan} \\
OE Secondary Nasalization & \emph{*líornōjąn} \\
OE I Umlaut & \emph{*líornējąn} \\
OE Unstressed Long Vowel Shortening & \emph{*líornejąn} \\
OE Weak Tail Reduction & \emph{*líornejan} \\
OE Intervocalic J Vocalization & \emph{*líorneian} \\
OE Unstressed EI Contraction & \emph{*líornian} \\
\end{tabular}
\end{minipage}
\\
\end{tabularx}
\end{minipage}%
} \endgroup

Ringe and Taylor's class-II weak family and Kroonen's comparative base
both support the same inherited verb, so the question here is not lexeme
choice but dialectal targeting (\citeproc{ref-RingeTaylor2014}{Ringe and
Taylor 2014, 318}; \citeproc{ref-Kroonen2013}{Kroonen 2013, 328}).
Campbell's discussion is decisive for that point: Northumbrian
\emph{liornian} preserves \emph{io}, whereas the more familiar
\emph{leornian} reflects a later \emph{eo} development
(\citeproc{ref-Campbell1959}{Campbell 1959, 123};
\citeproc{ref-Campbell1959}{1959, 296}). The Old English evidence
therefore does not force the dictionary headword. It allows the book to
keep the Northumbrian form as the regular comparator for this branch.

Dialect note. The trace shows the regular phonological route from
\emph{*líznōjaną} to \emph{liornian}; the real interpretive burden is
simply that most reference works privilege \emph{leornian} instead
(\citeproc{ref-ClarkHall1960}{Clark Hall 1960, 193};
\citeproc{ref-BrightCassidyRingler1971}{Bright 1971, 58}). In fuller
treatment, the Northumbrian and West Saxon contrast may deserve a
dedicated editorial comparison, but the essential point is already
contained in the Northumbrian evidence itself.

\subsection{linden --- OE lind}\label{linden-oe-lind}

Derivation: \emph{*líndō} \textgreater{} \emph{lind} (regular).

\begingroup
\setlength{\fboxsep}{6pt}

\noindent\fbox{%
\begin{minipage}{0.97\linewidth}
\small
\begin{tabularx}{\linewidth}{@{}>{\raggedright\arraybackslash}p{0.460\linewidth}>{\centering\arraybackslash}X>{\raggedright\arraybackslash}p{0.460\linewidth}@{}}
\begin{minipage}[t]{\linewidth}
\centering\textbf{Earlier Germanic changes}\par
\vspace{0.35em}
\raggedright
\centering\textbf{West Germanic}\par
\raggedright
\vspace{0.2em}
\raggedright [no change]\par
\vspace{0.6em}
\centering\textbf{Northwest Germanic}\par
\raggedright
\vspace{0.2em}
\begin{tabular}{@{}>{\raggedright\arraybackslash}p{0.74\linewidth}@{\hspace{0.55em}}>{\raggedright\arraybackslash}p{0.16\linewidth}@{\hspace{0.25em}}}
\mbox{NWGmc Final Long O Raising} & \emph{*líndu} \\
\end{tabular}
\end{minipage}
&
&
\begin{minipage}[t]{\linewidth}
\centering\textbf{Old English changes}\par
\vspace{0.35em}
\raggedright
\begin{tabular}{@{}>{\raggedright\arraybackslash}p{0.74\linewidth}@{\hspace{0.55em}}>{\raggedright\arraybackslash}p{0.16\linewidth}@{\hspace{0.25em}}}
\mbox{OE High Vowel Apocope} & \emph{*línd} \\
\end{tabular}
\end{minipage}
\\
\end{tabularx}
\end{minipage}%
} \endgroup

Kroonen's \emph{*lindō-} and the Old English dictionary evidence both
point to the ordinary noun \emph{lind} `linden, lime-tree'
(\citeproc{ref-Kroonen2013}{Kroonen 2013, 337};
\citeproc{ref-ClarkHall1960}{Clark Hall 1960, 198};
\citeproc{ref-BosworthToller1898}{Bosworth and Toller 1898, 630}). The
derivation is correspondingly brief: Northwest Germanic final \emph{ō}
raises to a \emph{*líndu} stage, and high-vowel apocope gives
\emph{lind}. No further explanation is needed beyond noting that Clark
Hall's adjectival \emph{linden} `made of linden-wood' belongs to a
different lexical use, not to the noun chosen here for direct comparison
(\citeproc{ref-ClarkHall1960}{Clark Hall 1960, 198}).

\subsection{mother --- OE mōder}\label{mother-oe-mux14dder}

Derivation: \emph{*mōdēr} \textgreater{} \emph{mōder} (regular).

\begingroup
\setlength{\fboxsep}{6pt}

\noindent\fbox{%
\begin{minipage}{0.97\linewidth}
\small
\begin{tabularx}{\linewidth}{@{}>{\raggedright\arraybackslash}p{0.460\linewidth}>{\centering\arraybackslash}X>{\raggedright\arraybackslash}p{0.460\linewidth}@{}}
\begin{minipage}[t]{\linewidth}
\centering\textbf{Earlier Germanic changes}\par
\vspace{0.35em}
\raggedright
\centering\textbf{West Germanic}\par
\raggedright
\vspace{0.2em}
\raggedright [no change]\par
\vspace{0.6em}
\centering\textbf{Northwest Germanic}\par
\raggedright
\vspace{0.2em}
\begin{tabular}{@{}>{\raggedright\arraybackslash}p{0.74\linewidth}@{\hspace{0.55em}}>{\raggedright\arraybackslash}p{0.16\linewidth}@{\hspace{0.25em}}}
\mbox{NWGmc Long E Lowering} & \emph{*mōdǣr} \\
\end{tabular}
\end{minipage}
&
&
\begin{minipage}[t]{\linewidth}
\centering\textbf{Old English changes}\par
\vspace{0.35em}
\raggedright
\begin{tabular}{@{}>{\raggedright\arraybackslash}p{0.74\linewidth}@{\hspace{0.55em}}>{\raggedright\arraybackslash}p{0.16\linewidth}@{\hspace{0.25em}}}
OE Unstressed Long Vowel Shortening & \emph{*mōdær} \\
OE Unstressed AE Merger & \emph{*mōder} \\
\end{tabular}
\end{minipage}
\\
\end{tabularx}
\end{minipage}%
} \endgroup

The compact issue here is not whether the lexeme is right but which Old
English form best represents the regular nominative line. Kroonen and
Orel give the Proto-Germanic r-stem as \emph{*mōdēr-} or \emph{*mōdēr},
while Old English reference works more often print \emph{mōdor} or
\emph{modor} and show oblique \emph{mēder} elsewhere in the paradigm
(\citeproc{ref-Kroonen2013}{Kroonen 2013, 368};
\citeproc{ref-Orel2003}{Orel 2003, 274};
\citeproc{ref-ClarkHall1960}{Clark Hall 1960, 219};
\citeproc{ref-Campbell1959}{Campbell 1959, 598};
\citeproc{ref-RingeTaylor2014}{Ringe and Taylor 2014, 312}). The trace
therefore supports a regular nominative \emph{mōder}, not the later
levelled dictionary headword. That is the point of compression in this
entry: the philology lies in paradigm history, not in an irregular sound
change.

Comparison note. The familiar \emph{mōdor} should remain visible in
fuller editorial treatment, because readers will expect it, but it does
not replace the regular nominative reflex represented here.

\subsection{show --- OE sċēawian}\label{show-oe-sux10bux113awian}

Derivation: \emph{*skáwōjaną} \textgreater{} \emph{sċēawian} (regular).

\begingroup
\setlength{\fboxsep}{6pt}

\noindent\fbox{%
\begin{minipage}{0.97\linewidth}
\footnotesize
\begin{tabularx}{\linewidth}{@{}>{\raggedright\arraybackslash}p{0.280\linewidth}>{\centering\arraybackslash}X>{\raggedright\arraybackslash}p{0.560\linewidth}@{}}
\begin{minipage}[t]{\linewidth}
\centering\textbf{Earlier Germanic changes}\par
\vspace{0.35em}
\raggedright
\centering\textbf{West Germanic}\par
\raggedright
\vspace{0.2em}
\raggedright [no change]\par
\vspace{0.6em}
\centering\textbf{Northwest Germanic}\par
\raggedright
\vspace{0.2em}
\raggedright [no change]\par
\end{minipage}
&
&
\begin{minipage}[t]{\linewidth}
\centering\textbf{Old English changes}\par
\vspace{0.35em}
\raggedright
\begin{tabular}{@{}>{\raggedright\arraybackslash}p{0.62\linewidth}@{\hspace{0.55em}}>{\raggedright\arraybackslash}p{0.28\linewidth}@{\hspace{0.25em}}}
OE Aw Long Diphthong & \emph{*skḗawōjaną} \\
OE Heavy Syllable Nasal Apocope & \emph{*skḗawōjan} \\
OE Secondary Nasalization & \emph{*skḗawōjąn} \\
OE Sk Palatalization & \emph{*ʃḗawōjąn} \\
OE I Umlaut & \emph{*ʃḗawējąn} \\
OE Unstressed Long Vowel Shortening & \emph{*ʃḗawejąn} \\
OE Weak Tail Reduction & \emph{*ʃḗawejan} \\
OE Intervocalic J Vocalization & \emph{*ʃḗaweian} \\
OE Unstressed EI Contraction & \emph{*ʃḗawian} \\
\end{tabular}
\end{minipage}
\\
\end{tabularx}
\end{minipage}%
} \endgroup

Orel and Kroonen both place the verb in the ordinary class-II
show-family, and Brunner's Old English evidence \emph{scēawian},
\emph{scāwian} confirms that the entry is dealing with the standard
lexeme rather than a special paradigm cell (\citeproc{ref-Orel2003}{Orel
2003, 339}; \citeproc{ref-Kroonen2013}{Kroonen 2013, 482};
\citeproc{ref-SieversBrunner1965}{Sievers 1965, 148}). The regularity
lies in the treatment of \emph{aw} before a following vowel: the
inherited form develops to Old English \emph{ēaw}, while the class-II
suffix preserves \emph{ō} between \emph{w} and \emph{j}, yielding the
expected \emph{sċēawian} (\citeproc{ref-Campbell1959}{Campbell 1959,
138}; \citeproc{ref-Orel2003}{Orel 2003, 339}). The remaining wrinkle is
orthographic, not philological.

Form note. The difference between dictionary \emph{scēawian} and the
book spelling \emph{sċēawian} is only normalization of initial
(\citeproc{ref-Campbell1959}{Campbell 1959, 19};
\citeproc{ref-Hogg1992}{Hogg 1992, 236}).

\subsection{weapon --- OE wǣpn}\label{weapon-oe-wux1e3pn}

Derivation: \emph{*wḗpną} \textgreater{} \emph{wǣpn} (regular).

\begingroup
\setlength{\fboxsep}{6pt}

\noindent\fbox{%
\begin{minipage}{0.97\linewidth}
\small
\begin{tabularx}{\linewidth}{@{}>{\raggedright\arraybackslash}p{0.460\linewidth}>{\centering\arraybackslash}X>{\raggedright\arraybackslash}p{0.460\linewidth}@{}}
\begin{minipage}[t]{\linewidth}
\centering\textbf{Earlier Germanic changes}\par
\vspace{0.35em}
\raggedright
\centering\textbf{West Germanic}\par
\raggedright
\vspace{0.2em}
\raggedright [no change]\par
\vspace{0.6em}
\centering\textbf{Northwest Germanic}\par
\raggedright
\vspace{0.2em}
\begin{tabular}{@{}>{\raggedright\arraybackslash}p{0.74\linewidth}@{\hspace{0.55em}}>{\raggedright\arraybackslash}p{0.16\linewidth}@{\hspace{0.25em}}}
\mbox{NWGmc Long E Lowering} & \emph{*wǣpną} \\
\end{tabular}
\end{minipage}
&
&
\begin{minipage}[t]{\linewidth}
\centering\textbf{Old English changes}\par
\vspace{0.35em}
\raggedright
\begin{tabular}{@{}>{\raggedright\arraybackslash}p{0.74\linewidth}@{\hspace{0.55em}}>{\raggedright\arraybackslash}p{0.16\linewidth}@{\hspace{0.25em}}}
OE Heavy Syllable Nasal Apocope & \emph{*wǣpn} \\
\end{tabular}
\end{minipage}
\\
\end{tabularx}
\end{minipage}%
} \endgroup

Kroonen's double-stem noun \emph{*wēbna- \textasciitilde{} wēpna-} and
Campbell's discussion of cluster nouns both show why this entry cannot
be reduced to the familiar late West Saxon headword \emph{wǣpen}
(\citeproc{ref-Kroonen2013}{Kroonen 2013, 617};
\citeproc{ref-Campbell1959}{Campbell 1959, 150};
\citeproc{ref-Campbell1959}{1959, 226}). Old English evidence preserves
both broken and unbroken forms: Bright contrasts broken nominative
\emph{wǣpen/wapen} with unbroken \emph{wǣpnes}, and Clark Hall likewise
records the headword under \emph{wapen} while preserving unbroken forms
elsewhere (\citeproc{ref-BrightCassidyRingler1971}{Bright 1971, 29};
\citeproc{ref-ClarkHall1960}{Clark Hall 1960, 355}). The point of the
book form \emph{wǣpn} is to represent the regular unbroken cluster
word-finally, which the trace derives cleanly from \emph{*wḗpną}.

Form note. The common simplex headword \emph{wǣpen} is real Old English,
but it reflects later breaking in the simplex. The unbroken forms remain
essential evidence for the regular line modeled here
(\citeproc{ref-Campbell1959}{Campbell 1959, 150};
\citeproc{ref-Campbell1959}{1959, 226};
\citeproc{ref-BrightCassidyRingler1971}{Bright 1971, 29};
\citeproc{ref-ClarkHall1960}{Clark Hall 1960, 355}).

\clearpage

\subsection*{References}\label{references}
\addcontentsline{toc}{subsection}{References}

\protect\phantomsection\label{refs}
\begin{CSLReferences}{1}{1}
\bibitem[\citeproctext]{ref-BosworthToller1898}
Bosworth, Joseph, and T. Northcote Toller. 1898. \emph{An {Anglo-Saxon}
Dictionary}. Clarendon Press.

\bibitem[\citeproctext]{ref-BrightCassidyRingler1971}
Bright, James W. 1971. \emph{Bright's {Old English} Grammar and Reader}.
3rd edn. Edited by Frederic G. Cassidy and Richard N. Ringler. Holt,
Rinehart; Winston.

\bibitem[\citeproctext]{ref-Campbell1959}
Campbell, Alistair. 1959. \emph{Old {English} Grammar}. Clarendon Press.

\bibitem[\citeproctext]{ref-ClarkHall1960}
Clark Hall, J. R. 1960. \emph{A Concise {Anglo-Saxon} Dictionary}. 4th
edn. University of Toronto Press.

\bibitem[\citeproctext]{ref-Fulk2018}
Fulk, R. D. 2018. \emph{A Comparative Grammar of the Early {Germanic}
Languages}. Vol. 3. Studies in Germanic Linguistics. John Benjamins.

\bibitem[\citeproctext]{ref-Hogg1992}
Hogg, Richard M. 1992. \emph{A Grammar of Old English. Volume 1:
Phonology}. Blackwell.

\bibitem[\citeproctext]{ref-Kroonen2013}
Kroonen, Guus. 2013. \emph{Etymological Dictionary of {Proto-Germanic}}.
Vol. 11. Leiden Indo-European Etymological Dictionary Series. Brill.

\bibitem[\citeproctext]{ref-Orel2003}
Orel, Vladimir. 2003. \emph{A Handbook of {Germanic} Etymology}. Brill.

\bibitem[\citeproctext]{ref-RingeTaylor2014}
Ringe, Don, and Ann Taylor. 2014. \emph{The Development of {Old
English}}. Vol. 2. A Linguistic History of English. Oxford University
Press.

\bibitem[\citeproctext]{ref-SieversBrunner1965}
Sievers, Eduard. 1965. \emph{Altenglische Grammatik}. 3rd edn. Edited by
Karl Brunner. Niemeyer.

\end{CSLReferences}

\end{document}
